\begin{abstract}
Late-interaction multilingual retrievers such as ColBERT score query-document pairs at the token level via MaxSim, placing the geometry of individual token vectors at the center of retrieval quality. We show that multilingual ColBERT trained with bitext InfoNCE produces token vectors that are anisotropic: they concentrate in a narrow cone with effective rank well below the embedding dimension, with the target-side language collapsing harder than the source-side. To address this, we introduce \textbf{IsoColBERT}, a single regularizer term added to the InfoNCE training objective that penalizes the mean squared off-diagonal cosine similarity of valid token vectors in each batch, with zero cost at inference. Across three random seeds on FLORES+, IsoColBERT improves MaxSim P@1 by $3.45$ percentage points on English-Swahili (every seed individually positive) and $1.99$ percentage points on English-Arabic, with the gain extending zero-shot to two further low-resource languages (Nepali $+1.46$ pp, Tamil $+0.88$ pp), while leaving high-resource pairs unchanged. Mechanism ablations rule out localized and matching-level placements as substitutes for full-manifold regularization. The gains concentrate where pretrained signal is weakest, suggesting that token geometry is a tractable lever for low-resource multilingual retrieval.
\end{abstract}
